\chapter{Generated Examples and Probability Overlays}
\makeatletter
% In-place (non-floating) figures and tables: each example stays beside the
% question it answers.
\newenvironment{herefigure}{\par\medskip\noindent\begin{minipage}{\linewidth}\centering\def\@captype{figure}}{\end{minipage}\par\medskip}
\newenvironment{heretable}{\par\medskip\noindent\begin{minipage}{\linewidth}\centering\def\@captype{table}}{\end{minipage}\par\medskip}
\makeatother
\label{ch:generated-examples}

\index{generated examples}
\index{probability overlay}
Most specifications explain themselves with hand-written examples: a trace
someone typed, a sequence diagram someone drew. Those examples drift. This
chapter takes the opposite approach. Every example below is the answer to a
question asked of the specification, and every number is computed from it.
Nothing in the figures was drawn by hand; regenerate them with
\texttt{lake env lean --run ProbabilityOverlayExample.lean}.

The reading of the temporal operators is generative. If we generate event
lists from the grammar, then
\begin{description}
\item[$EF\,p$] means some generated list contains $p$; the witness is that
list, and it is the example.
\item[$AG\,p$] means every generated list satisfies $p$; a counterexample is a
generated list that does not.
\item[$AF\,p$] means every generated list eventually reaches $p$; a
counterexample is a list that loops forever without it.
\end{description}

\section{The core says what can happen}

A client sends a \texttt{Request}. The server answers with a
\texttt{Response}, or a \texttt{Timer} fires a \texttt{Timeout} and the client
retries. After the third timeout the client gives up. Retries are bounded, so
the task unrolls them into states: one waiting state per attempt.

The core carries no numbers. It says only which messages can follow which, and
it already answers the qualitative questions: every run ends, and giving up is
possible.

\section{The overlay says how likely and how long}

A separate table, the \emph{overlay}, annotates each transition with a branch
probability and a delivery delay. This one is \emph{declared}: it states what
we expect (a 10\% timeout rate, 40\,ms responses, 200\,ms timeouts). A
\emph{measured} overlay would have exactly the same shape, reduced from
recorded traces instead of asserted, and the same questions could be asked of
either. Keeping the numbers out of the grammar means a project without
probabilities loses nothing, and declared and measured values can be compared
on one unchanged specification.

The overlay is validated before use: every transition must be annotated,
nothing else may be, and the probabilities leaving each state must sum to
exactly~1. They are exact fractions and are never normalized silently.

\begin{herefigure}
\AccessibleFigure{State machine of the bounded retry task with the declared overlay on its edges. idle sends Request with certainty after 5 ms to waiting1. Each waiting state receives Response with probability 9/10 after 40 ms and moves to done or receives Timeout with probability 1/10 after 200 ms. The first two timeouts lead to a backoff state that sends Request again after 5 ms. The third timeout leads to gave up.}{\resizebox{!}{.55\textheight}{% Generated by ProbabilityOverlayExample.lean from LeanFM.Overlay.declared.
\begin{tikzpicture}[>=Stealth,state/.style={draw=diagramink,text=diagramink,fill=diagramsurface,rounded corners,minimum width=1.7cm,minimum height=7mm,font=\small},edge/.style={->,draw=leanblue,thick},lbl/.style={font=\scriptsize,text=diagramink,align=center,fill=pagebackground,inner sep=1pt}]
\node[state] (idle) at (0,0) {idle};
\node[state] (waiting1) at (0,-1.7) {waiting1};
\node[state] (backoff1) at (0,-3.4) {backoff1};
\node[state] (waiting2) at (0,-5.1) {waiting2};
\node[state] (backoff2) at (0,-6.8) {backoff2};
\node[state] (waiting3) at (0,-8.5) {waiting3};
\node[state,double] (done) at (6,-5.1) {done};
\node[state,double] (gaveup) at (0,-10.2) {gave\_up};
\draw[edge] (idle) to[] node[lbl,left=3pt] {Request\\ 5ms} (waiting1);
\draw[edge] (waiting1) to[bend left=10] node[lbl] {Response\\ p=9/10, 40ms} (done);
\draw[edge] (waiting1) to[] node[lbl,left=3pt] {Timeout\\ p=1/10, 200ms} (backoff1);
\draw[edge] (backoff1) to[] node[lbl,left=3pt] {Request\\ 5ms} (waiting2);
\draw[edge] (waiting2) to[] node[lbl] {Response\\ p=9/10, 40ms} (done);
\draw[edge] (waiting2) to[] node[lbl,left=3pt] {Timeout\\ p=1/10, 200ms} (backoff2);
\draw[edge] (backoff2) to[] node[lbl,left=3pt] {Request\\ 5ms} (waiting3);
\draw[edge] (waiting3) to[bend right=10] node[lbl] {Response\\ p=9/10, 40ms} (done);
\draw[edge] (waiting3) to[] node[lbl,left=3pt] {Timeout\\ p=1/10, 200ms} (gaveup);
\end{tikzpicture}
}}
\caption{The bounded-retry task (the core) with the declared overlay drawn on
its edges. Double borders mark terminal states.}
\label{fig:overlay-state-machine}
\end{herefigure}

\section{Questions and their generated answers}

\begin{heretable}
\small
\begin{tabular}{p{.3\linewidth}p{.3\linewidth}p{.32\linewidth}}
\textbf{Question} & \textbf{Formula} & \textbf{Answer}\\
\hline
% Generated by ProbabilityOverlayExample.lean.
Does every run end? (core, no overlay) & \texttt{AF terminal} & proved \\
Can the client give up? & \texttt{EF gave\_up} & proved; most likely witness has p=1/1000, 615 ms \\
How likely is success? & \texttt{P=? [F done]} & 999/1000 (0.999) \\
Do 99\% succeed within 300 ms? & \texttt{P>=99/100 [F<=300ms done]} & proved; exactly 99/100 \\
Do 99.9\% succeed within 300 ms? & \texttt{P>=999/1000 [F<=300ms done]} & disproved; most likely late run p=9/1000, 455 ms \\
Expected duration of a run & \texttt{E[duration]} & 67.71 ms \\

\end{tabular}
\caption{Every answer is computed from the task and the declared overlay. The
first row uses the core alone.}
\label{tab:overlay-answers}
\end{heretable}

The table reads like a service-level review. Success is overwhelmingly likely
($999/1000$). Ninety-nine percent of runs finish within 300\,ms, and that
bound is met \emph{exactly}, with no margin. Tightening it to 99.9\% fails, and
the checker does not just say so: it produces the most likely run that misses
the deadline.

\subsection{The usual run}

The most likely generated run is the one a newcomer should see first.

\begin{herefigure}
\AccessibleFigure{Most likely generated run. Message from Client to Server at 5 ms Request. Message from Server to Client at 45 ms Response with probability 9/10.}{\resizebox{.7\textwidth}{!}{% Generated by LeanFM.MDPDiagram.renderInteractionTeX.
\begin{tikzpicture}[actor/.style={draw=diagramink,text=diagramink,fill=diagramsurface,rounded corners,minimum width=2.3cm,minimum height=7mm,font=\small},message/.style={->,>=Stealth,draw=leanblue,thick}]
\node[actor] (a0) at (0,0) {Client};
\draw[dashed,draw=diagramink] (a0.south) -- ++(0,-3);
\node[actor] (a1) at (6,0) {Server};
\draw[dashed,draw=diagramink] (a1.south) -- ++(0,-3);
\node[actor] (a2) at (12,0) {Timer};
\draw[dashed,draw=diagramink] (a2.south) -- ++(0,-3);
\draw[message] (0,-1) -- node[above,fill=pagebackground,text=diagramink,font=\scriptsize] {t=5ms Request} (6,-1);
\draw[message] (6,-2) -- node[above,fill=pagebackground,text=diagramink,font=\scriptsize] {t=45ms Response (p=9/10)} (0,-2);
\end{tikzpicture}
}}
\caption{The most likely run: one request and one response, 45\,ms, probability
$9/10$.}
\label{fig:overlay-run-usual}
\end{herefigure}

\subsection{A precedent: the client gives up}

$EF\,\mathit{gave\_up}$ holds, and its witness is a concrete case. It is rare,
with probability $1/1000$, but it is possible, and this is what it looks like.

\begin{herefigure}
\AccessibleFigure{Generated witness for EF gave up. Three rounds of Request from Client to Server each followed by Timeout from Timer to Client with probability 1/10. The final timeout arrives at 615 ms.}{\resizebox{.85\textwidth}{!}{% Generated by LeanFM.MDPDiagram.renderInteractionTeX.
\begin{tikzpicture}[actor/.style={draw=diagramink,text=diagramink,fill=diagramsurface,rounded corners,minimum width=2.3cm,minimum height=7mm,font=\small},message/.style={->,>=Stealth,draw=leanblue,thick}]
\node[actor] (a0) at (0,0) {Client};
\draw[dashed,draw=diagramink] (a0.south) -- ++(0,-7);
\node[actor] (a1) at (6,0) {Server};
\draw[dashed,draw=diagramink] (a1.south) -- ++(0,-7);
\node[actor] (a2) at (12,0) {Timer};
\draw[dashed,draw=diagramink] (a2.south) -- ++(0,-7);
\draw[message] (0,-1) -- node[above,fill=pagebackground,text=diagramink,font=\scriptsize] {t=5ms Request} (6,-1);
\draw[message] (12,-2) -- node[above,fill=pagebackground,text=diagramink,font=\scriptsize] {t=205ms Timeout (p=1/10)} (0,-2);
\draw[message] (0,-3) -- node[above,fill=pagebackground,text=diagramink,font=\scriptsize] {t=210ms Request} (6,-3);
\draw[message] (12,-4) -- node[above,fill=pagebackground,text=diagramink,font=\scriptsize] {t=410ms Timeout (p=1/10)} (0,-4);
\draw[message] (0,-5) -- node[above,fill=pagebackground,text=diagramink,font=\scriptsize] {t=415ms Request} (6,-5);
\draw[message] (12,-6) -- node[above,fill=pagebackground,text=diagramink,font=\scriptsize] {t=615ms Timeout (p=1/10)} (0,-6);
\end{tikzpicture}
}}
\caption{The witness for $EF\,\mathit{gave\_up}$: three timeouts, 615\,ms,
probability $1/1000$.}
\label{fig:overlay-run-giveup}
\end{herefigure}

\subsection{A counterexample: the late success}

The claim $P_{\ge 999/1000}[F_{\le 300\,\mathrm{ms}}\,\mathit{done}]$ fails. Of
all the runs that miss the deadline, the most likely is not the give-up run
but a success on the third attempt. That is the run worth fixing first.

\begin{herefigure}
\AccessibleFigure{Most likely run that misses the 300 ms deadline. Two rounds of Request then Timeout with probability 1/10 each followed by a third Request and a Response with probability 9/10 at 455 ms.}{\resizebox{.85\textwidth}{!}{% Generated by LeanFM.MDPDiagram.renderInteractionTeX.
\begin{tikzpicture}[actor/.style={draw=diagramink,text=diagramink,fill=diagramsurface,rounded corners,minimum width=2.3cm,minimum height=7mm,font=\small},message/.style={->,>=Stealth,draw=leanblue,thick}]
\node[actor] (a0) at (0,0) {Client};
\draw[dashed,draw=diagramink] (a0.south) -- ++(0,-7);
\node[actor] (a1) at (6,0) {Server};
\draw[dashed,draw=diagramink] (a1.south) -- ++(0,-7);
\node[actor] (a2) at (12,0) {Timer};
\draw[dashed,draw=diagramink] (a2.south) -- ++(0,-7);
\draw[message] (0,-1) -- node[above,fill=pagebackground,text=diagramink,font=\scriptsize] {t=5ms Request} (6,-1);
\draw[message] (12,-2) -- node[above,fill=pagebackground,text=diagramink,font=\scriptsize] {t=205ms Timeout (p=1/10)} (0,-2);
\draw[message] (0,-3) -- node[above,fill=pagebackground,text=diagramink,font=\scriptsize] {t=210ms Request} (6,-3);
\draw[message] (12,-4) -- node[above,fill=pagebackground,text=diagramink,font=\scriptsize] {t=410ms Timeout (p=1/10)} (0,-4);
\draw[message] (0,-5) -- node[above,fill=pagebackground,text=diagramink,font=\scriptsize] {t=415ms Request} (6,-5);
\draw[message] (6,-6) -- node[above,fill=pagebackground,text=diagramink,font=\scriptsize] {t=455ms Response (p=9/10)} (0,-6);
\end{tikzpicture}
}}
\caption{The most likely counterexample to the 99.9\% / 300\,ms objective:
success on the third attempt at 455\,ms, probability $9/1000$.}
\label{fig:overlay-run-late}
\end{herefigure}

\section{The distribution behind the numbers}

The completion-time distribution is computed exactly by enumerating the
generated runs; it is not sampled. The logarithmic scale shows the long tail
that a mean of 67.71\,ms hides.

\begin{herefigure}
\AccessibleFigure{Bar chart of completion time for runs that end in done on a logarithmic probability axis. 45 ms has probability 9/10. 250 ms has probability 9/100. 455 ms has probability 9/1000.}{% Generated by ProbabilityOverlayExample.lean: completion-time distribution of runs ending in done.
\begin{tikzpicture}
\begin{axis}[ybar,ymode=log,log origin=infty,width=.92\linewidth,height=5.2cm,bar width=16pt,xlabel={completion time (ms)},ylabel={probability},symbolic x coords={45,250,455},xtick=data,nodes near coords,nodes near coords style={font=\scriptsize},point meta=explicit symbolic,ymin=0.0005,ymax=3,enlarge x limits=0.18]
\addplot[fill=leanblue,draw=leanblue] coordinates {(45,0.900000) [9/10](250,0.090000) [9/100](455,0.009000) [9/1000]};
\end{axis}
\end{tikzpicture}
}
\caption{Completion times of successful runs under the declared overlay. The
remaining $1/1000$ gives up at 615\,ms.}
\label{fig:overlay-distribution}
\end{herefigure}

\section{What comes next}

The same functions accept a measured overlay: count how often each branch was
taken in recorded traces and measure the time between each event and its
cause. Asking the questions in Table~\ref{tab:overlay-answers} of both overlays
turns ``we expected a 10\% timeout rate'' into a checked comparison, and the
branch where they disagree is the one to highlight. Retry loops here are
unrolled into bounded states; unbounded loops need the linear-equation
solution of probabilistic model checking rather than run enumeration, and the
current implementation reports them instead of truncating them.
